

\documentclass[10pt, oneside]{article}   	% use "amsart" instead of "article" for AMSLaTeX format
\usepackage{geometry}                		% See geometry.pdf to learn the layout options. There are lots.
\geometry{a4paper}                   		% ... or a4paper or a5paper or ... 
%\geometry{landscape}                		% Activate for rotated page geometry
%\usepackage[parfill]{parskip}    		% Activate to begin paragraphs with an empty line rather than an indent
\usepackage{graphicx}				% Use pdf, png, jpg, or eps§ with pdflatex; use eps in DVI mode
								% TeX will automatically convert eps --> pdf in pdflatex		
\usepackage{amssymb}
\usepackage{amsmath}
\usepackage{authblk}
\usepackage[hyphens]{url}

\usepackage{url}
\usepackage{hyperref} 

\usepackage{titlesec}

\usepackage{caption}

\usepackage{float}


\title{A Comprehensive Approach to Ecosystem Services Management: From Alarm to Action}
\author{Aidan T. Parkinson and Andreea M. Benu, 5 Colina Mews, London, N15 3HS}

\date{\today}							% Activate to display a given date or no date

\begin{document}
\maketitle
\begin{center}
\href{mailto:enquiries@aidantparkinson.com}{enquiries@aidantparkinson.com}
\end{center}

\section{Abstract}
This paper examines how to develop a reasonable contractual proposition for a society to cooperate fairly in the interest of a common ecosystem.
It recognises that no society is entirely free from political realities or the costs involved in maintaining social order.
As a result, ideal propositions of social cooperation, such as John Rawls’s The Law of Peoples, remain incomplete when considering practical application.
We introduce the Commonwealth Cost of Carbon, as a new Ecosystem Performance Observation that measures how much societies spend to maintain order relative to their total activity.
This ratio provides a transparent and replicable way to assess global ecosystem performance without relying on expert computing resources or privileged information.
Over just twenty years, this cost increased from $\sim$2000\$65 to $\sim$2020\$140 per tonne of CO2e, suggesting both a decline in ecosystem performance and a rising social value for emission reduction.
Building on this indicator, we propose the Commonwealth of Peoples as a new trade association, that translates moral duty into service performance.
Here, a complementary Service Rating is introduced, serving as a cost–benefit ratio that benchmarks how industrialised is a habitat.
Together, these tools provide a performance-based foundation of ecosystem governance that provides an assured approach to move a society from alarm to action.\\

Keywords: Ecosystem Performance; Ecological Economics; Social Contract; Trade Association; Service Performance Marketing\\

\end{document}